% VI49-DETAILED-EXERCISE-SOLUTIONS
% VI49-DETAILED-EXERCISE-01
\subsection*{Exercise VI.49.1}
\begin{solution}
For an open \(U\subseteq X\), every restriction map of \(\mathcal F|_U\) is already a restriction map of \(\mathcal F\).  Thus a flabby sheaf remains flabby after restriction.  Flabby acyclicity from VI/46 then gives
\[
H^q(U,\mathcal F|_U)=0\qquad(q>0).
\]
\end{solution}

% VI49-DETAILED-EXERCISE-02
\subsection*{Exercise VI.49.2}
\begin{solution}
A skyscraper sheaf is flabby: sections over an open are either the fiber group or zero, and every restriction is surjective.  Hence flabby vanishing gives zero higher cohomology.  Finite sums of skyscrapers are therefore acyclic as well.
\end{solution}

% VI49-DETAILED-EXERCISE-03
\subsection*{Exercise VI.49.3}
\begin{solution}
If \(N_{f_i}=0\), then for each \(n\in N\) some power \(f_i^{m_i}\) annihilates \(n\).  Choose a common large \(m\).  Since the \(f_i\) generate the unit ideal, so do sufficiently large powers after adjusting coefficients: \(1=\sum a_i f_i^m\).  Multiplying by \(n\) gives \(n=0\), hence \(N=0\).
\end{solution}

% VI49-DETAILED-EXERCISE-04
\subsection*{Exercise VI.49.4}
\begin{solution}
For \(D(f)\cup D(g)=\operatorname{Spec}A\), the augmented localization complex is
\[
0\to M\longrightarrow M_f\oplus M_g\longrightarrow M_{fg}\to0,
\]
where the second map is the difference of the two localization maps.  There are no higher terms because the cover has only two members.
\end{solution}

% VI49-DETAILED-EXERCISE-05
\subsection*{Exercise VI.49.5}
\begin{solution}
Let \(H\) be a homology module of the localization complex.  Localizing at \(f\) makes \(D(f)\) the whole localized space, so the augmented Čech complex contracts; hence \(H_f=0\).  Similarly \(H_g=0\).  Since \(D(f)\cup D(g)\) covers \(\operatorname{Spec}A\), a module vanishing after both localizations is zero, proving exactness.
\end{solution}

% VI49-DETAILED-EXERCISE-06
\subsection*{Exercise VI.49.6}
\begin{solution}
Affine vanishing applies directly:
\[
H^q(\operatorname{Spec}A,\widetilde M)=0\qquad(q>0).
\]
The degree-zero group is \(M\), so on an affine scheme quasi-coherent sheaf cohomology is completely concentrated in degree zero.
\end{solution}

% VI49-DETAILED-EXERCISE-07
\subsection*{Exercise VI.49.7}
\begin{solution}
A prime \(\mathfrak p\) lies in \(D(f_1)\cap\cdots\cap D(f_p)\) exactly when none of the \(f_i\) belongs to \(\mathfrak p\).  Since \(\mathfrak p\) is prime, this is equivalent to \(f_1\cdots f_p\notin\mathfrak p\).  Therefore the intersection equals \(D(f_1\cdots f_p)\).
\end{solution}

% VI49-DETAILED-EXERCISE-08
\subsection*{Exercise VI.49.8}
\begin{solution}
Every finite intersection of a principal cover is again a principal open \(D(f_{i_0}\cdots f_{i_p})\), hence affine.  Restricting a quasi-coherent sheaf remains quasi-coherent, and affine vanishing kills all positive cohomology on every intersection.  Thus the cover is Leray.
\end{solution}

% VI49-DETAILED-EXERCISE-09
\subsection*{Exercise VI.49.9}
\begin{solution}
For affine opens \(U,V\subset X\),
\[
U\cap V=\Delta_X^{-1}(U\times V).
\]
If \(X\) is separated, \(\Delta_X\) is a closed immersion.  The product \(U\times V\) is affine, and a closed subscheme of an affine scheme is affine.  Hence \(U\cap V\) is affine.
\end{solution}

% VI49-DETAILED-EXERCISE-10
\subsection*{Exercise VI.49.10}
\begin{solution}
A three-member affine cover of a separated scheme is Leray.  Its Čech complex has terms only in degrees \(0,1,2\).  Therefore
\[
H^q(X,\mathcal F)=0\qquad(q\ge3)
\]
for every quasi-coherent \(\mathcal F\).
\end{solution}

% VI49-DETAILED-EXERCISE-11
\subsection*{Exercise VI.49.11}
\begin{solution}
A two-member affine cover is Leray and its Čech complex stops in degree one.  Consequently
\[
H^q(X,\mathcal F)=0\qquad(q\ge2)
\]
for every quasi-coherent sheaf on the separated scheme.
\end{solution}

% VI49-DETAILED-EXERCISE-12
\subsection*{Exercise VI.49.12}
\begin{solution}
For \(n=2\), global sections are homogeneous quadrics:
\[
H^0(\mathbb P^1,\mathcal O(2))
=\langle x_0^2,x_0x_1,x_1^2\rangle_k.
\]
Thus the dimension is three.
\end{solution}

% VI49-DETAILED-EXERCISE-13
\subsection*{Exercise VI.49.13}
\begin{solution}
The standard compatibility calculation says \(H^0(\mathcal O(n))\) is spanned by Laurent exponents \(0,\ldots,n\).  For \(n=-1\) there are no such exponents, so
\[
H^0(\mathbb P^1,\mathcal O(-1))=0.
\]
\end{solution}

% VI49-DETAILED-EXERCISE-14
\subsection*{Exercise VI.49.14}
\begin{solution}
Using the Leray standard cover,
\[
H^1(\mathcal O(-1))
=\frac{k[t,t^{-1}]}{k[t]+t^{-1}k[t^{-1}]}.
\]
The two denominator summands contain all exponents \(\ge0\) and \(\le-1\), so the quotient is zero.
\end{solution}

% VI49-DETAILED-EXERCISE-15
\subsection*{Exercise VI.49.15}
\begin{solution}
For \(n=-4\), the denominator kills exponents \(\ge0\) and \(\le-4\).  The missing exponents are \(-3,-2,-1\), so
\[
H^1(\mathbb P^1,\mathcal O(-4))
=\langle t^{-3},t^{-2},t^{-1}\rangle_k.
\]
\end{solution}

% VI49-DETAILED-EXERCISE-16
\subsection*{Exercise VI.49.16}
\begin{solution}
For \(\mathcal O(-4)\), \(h^0=0\) and \(h^1=3\).  Hence
\[
\chi(\mathbb P^1,\mathcal O(-4))=0-3=-3,
\]
which agrees with the uniform formula \(\chi(\mathcal O(n))=n+1\).
\end{solution}

% VI49-DETAILED-EXERCISE-17
\subsection*{Exercise VI.49.17}
\begin{solution}
In a short exact sequence with \(\mathcal F'\) and \(\mathcal F\) acyclic, the first boundary has zero target, so global sections surject onto the quotient.  In positive degrees the long exact sequence places \(H^q(\mathcal F'')\) between zero groups.  Hence the quotient is acyclic.
\end{solution}

% VI49-DETAILED-EXERCISE-18
\subsection*{Exercise VI.49.18}
\begin{solution}
If the middle and quotient sheaves are acyclic, higher exactness already gives \(H^q(\mathcal F')=0\) for \(q\ge2\).  To kill \(H^1(\mathcal F')\) one must additionally require
\[
\Gamma(X,\mathcal F)\twoheadrightarrow\Gamma(X,\mathcal F'').
\]
Then the first connecting obstruction vanishes and the kernel is acyclic.
\end{solution}

% VI49-DETAILED-EXERCISE-19
\subsection*{Exercise VI.49.19}
\begin{solution}
From
\[
0\to\mathcal F'\to\mathcal I\to\mathcal F''\to0
\]
with \(\mathcal I\) acyclic, the long exact sequence contains
\[
0=H^q(\mathcal I)\to H^q(\mathcal F'')\to H^{q+1}(\mathcal F')\to H^{q+1}(\mathcal I)=0
\]
for \(q\ge1\).  Thus the middle arrow is a natural isomorphism.
\end{solution}

% VI49-DETAILED-EXERCISE-20
\subsection*{Exercise VI.49.20}
\begin{solution}
A coherent sheaf with finite support has finite length and admits a filtration whose successive quotients are skyscraper sheaves (possibly with finite-dimensional fibers).  Skyscraper sheaves are acyclic, and repeated use of the long exact sequence propagates acyclicity through the filtration.
\end{solution}

% VI49-DETAILED-EXERCISE-21
\subsection*{Exercise VI.49.21}
\begin{solution}
The point sequence
\[
0\to\mathcal O(n-1)\to\mathcal O(n)\to k_p\to0
\]
has an acyclic quotient.  The long exact sequence therefore gives a surjection
\[
H^1(\mathcal O(n-1))\twoheadrightarrow H^1(\mathcal O(n)),
\]
whose kernel is exactly the cokernel of the evaluation map \(H^0(\mathcal O(n))\to k_p\).
\end{solution}

% VI49-DETAILED-EXERCISE-22
\subsection*{Exercise VI.49.22}
\begin{solution}
Affine vanishing uses the equivalence between quasi-coherent sheaves and module-localization data on affine schemes.  An arbitrary sheaf of abelian groups or arbitrary \(\mathcal O_X\)-module need not arise from such localization data.  Therefore affineness alone does not make every sheaf acyclic.
\end{solution}

% VI49-DETAILED-EXERCISE-23
\subsection*{Exercise VI.49.23}
\begin{solution}
Separatedness guarantees that finite intersections of affine opens are affine, via the diagonal criterion.  Without separatedness an intersection may be nonaffine, so affine vanishing does not automatically apply on overlaps and the affine cover need not be Leray.
\end{solution}

% VI49-DETAILED-EXERCISE-24
\subsection*{Exercise VI.49.24}
\begin{solution}
Affine vanishing says that every quasi-coherent sheaf has zero positive cohomology on an affine scheme, with no positivity hypothesis.  Serre vanishing concerns coherent sheaves on projective schemes and states that sufficiently positive twists eventually kill higher cohomology.  It is asymptotic, global, and substantially deeper.
\end{solution}

